\section{Raw Representation Invention, Outcome Grounding, and Autonomous Curriculum}
\label{sec:aion-raw-representation-outcome-curriculum}

\subsection{Objective and Numbering Policy}

Following Phase~63, development is organized by named capability tracks rather
than overloaded numerical phase identifiers.  The objective of this programme
was to connect three learning layers:
\[
\begin{aligned}
\text{raw events}
&\rightarrow \text{invented representation}\\
&\rightarrow \text{independently observed outcomes}\\
&\rightarrow \text{autonomous curriculum}.
\end{aligned}
\]
No semantic facts such as \texttt{authority\_stale} or
\texttt{evidence\_conflict} were supplied to the new representation learner.

\subsection{Photon and GlyphNet Boundary}

The Photon documentation distinguishes two artifact families.  Files ending in
\texttt{.phn} or \texttt{.ptn} contain Photon Language and execute through the
Photon runtime.  Files ending in \texttt{.photon} or \texttt{.pthon} contain
glyph-compressed Python and execute through the CPython importer after
expansion.  The present work uses a data-only Photon Intermediate
Representation (PIR) for learned transition proposals:
\[
s_i \;\Rightarrow_{\alpha_j}\; s_k.
\]
Photon remains an executable proposal representation; HexCore verification and
CAU remain the final authority.

GlyphNet's append-only event ledger, source hashes, timestamps, stable
identifiers and graph edges are suitable raw inputs.  Existing semantic event
labels must be masked when testing vocabulary invention, otherwise the expected
ontology would leak into the learner.

\subsection{Raw Event Representation Learning}

The learner receives anonymous before/action/after records.  Each development
and sealed domain uses independent channel names, action names and categorical
value tokens.  It performs:
\begin{enumerate}
  \item anonymous channel--action effect-graph induction;
  \item iterative structural colour refinement without identifier strings;
  \item value orientation from interventions;
  \item predictive-equivalence state aggregation;
  \item action-conditioned relation induction;
  \item traceable Photon PIR compilation;
  \item structural transfer or explicit abstention.
\end{enumerate}

Symmetric dimensions remain a shared structural role; the system does not
invent unsupported semantic stories to distinguish them.  It invented five
predictive state symbols and thirteen transition symbols.

\begin{table}[ht]
\centering
\begin{tabular}{lr}
\toprule
Measure & Result \\
\midrule
Development families & 4 \\
Sealed renamed families & 4 \\
Semantic state/relation labels supplied & No \\
Shared raw field names & No \\
Mean and weakest state recovery & 100\% \\
Mean and weakest transition recovery & 100\% \\
Probes per transferred domain & 5 \\
Reduction versus exhaustive observation & 93.75\% \\
Non-isomorphic family & Abstain \\
Unsafe acceptances & 0 \\
Restart relearning & 0 \\
\bottomrule
\end{tabular}
\caption{Raw-event representation-learning result.}
\end{table}

The promoted procedure is
\texttt{procedure\_raw\_event\_representation\_24487fa43f20}.

\subsection{Real Outcome Grounding}

Four independent outcome authorities were attached:
\begin{itemize}
  \item safe AST execution for code candidates;
  \item exact disk rereading for documentary evidence;
  \item delayed withheld observations for forecasts;
  \item an operational precondition monitor for workflows.
\end{itemize}
Proposal confidence was never treated as truth.  The challenger improved from a
35\% ungrounded control to 100\% across twenty source-disjoint sealed tasks,
with 100\% weakest-family accuracy and zero unsafe acceptances.  The promoted
procedure is
\texttt{procedure\_real\_outcome\_grounding\_}\allowbreak
\texttt{1adb0a2eef1f}.

\subsection{Autonomous Capability Curriculum}

The curriculum controller computes a capability map from externally scored
accuracy and verification cost.  It selects the current weakest or
highest-cost family, generates a private curriculum, trains a reversible
specialist and audits backward retention before the next generation.

The autonomous order was delayed forecasting, document evidence, software
execution and operational execution.  All four generations were retained.

\begin{table}[ht]
\centering
\begin{tabular}{lrr}
\toprule
Measure & Control & Curriculum \\
\midrule
Source-disjoint arena tasks & 1000 & 1000 \\
Accuracy & 100\% & 100\% \\
Weakest-family accuracy & 100\% & 100\% \\
Mean independent verification probes & 2.486 & 1.000 \\
Probe reduction & --- & 59.77\% \\
Backward retention & --- & 100\% \\
Unsafe acceptances & --- & 0 \\
\bottomrule
\end{tabular}
\caption{Four-generation autonomous-curriculum arena.}
\end{table}

The promoted procedure is
\texttt{procedure\_autonomous\_curriculum\_24ae2ee1fbda}.
The capability map, specialists, sealed trace hash and complete generation
history survived restart with zero relearning.

\subsection{Claim Boundary}

This establishes a bounded chain from anonymous events to invented predictive
symbols, independently verified outcomes and autonomously selected curricula.
The event interface, finite action topology, task generators, descriptor
language and outcome authorities remain engineered.  It does not establish
unrestricted ontology induction, human-level commonsense or AGI.

The next decisive experiment is open-topology natural-event induction:
AION must infer entities, persistent and temporary properties, latent
dimensions and causal relations from unfamiliar natural project logs without
being told either the field vocabulary or the number of state variables.
